\section{光谱核对与预处理}
\subsection{原始数据的完整性}
四个附件均含一个工作表，首行为字段名称，随后各有7469组观测，共29876组。所有数值单元格可解析为有限数，波数严格递增，无重复波数，四份数据采用相同采样网格。波数范围为$399.6747$--$4000.1220\cmiv$，中位采样间隔为$0.4820\cmiv$。由于输入已经四舍五入，微小的间隔差异不能视为仪器扫描非均匀性的独立证据。

\begin{table}[htbp]\centering
\caption{四个附件的数据核对结果}\label{tab:quality}
\begin{tabular}{cccccc}
\toprule 附件 & 材料 & 角度 & 观测数 & 最大反射率/\% & 超100\%点数\\\midrule
1 & SiC & 10$^\circ$ & 7469 & 95.3850 & 0 \\
2 & SiC & 15$^\circ$ & 7469 & 102.7394 & 262 \\
3 & Si & 10$^\circ$ & 7469 & 79.8020 & 0 \\
4 & Si & 15$^\circ$ & 7469 & 91.4894 & 0 \\
\bottomrule
\end{tabular}\end{table}

\tabref{tab:quality}表明，四条光谱都在首个波数处出现一个反射率恰为零的点，而下一个采样点立即恢复到明显非零水平。仅凭附件无法确认它是仪器初始化标记还是其他端点处理结果。我们保留原始文件，不对其赋予物理含义；主拟合波段不包含该点，低波数比较也从下一个点开始。

附件2共有262点高于100\%，最大值为102.7394\%。超过理想绝对反射率范围说明参考归一化或标定值得检查，但并不能通过直接裁剪恢复真实值。裁剪会人为改变峰顶曲率、条纹振幅及谐波含量。因此，预处理不把大于100\%的数值压到100\%，而是将绝对增益与背景的不确定性留在模型中，并避免利用这些点倒推出未经标定的绝对光学常数。

\subsection{按材料选择分析波段}
\paperfigure{figures/raw_sic.png}{碳化硅双角度原始反射光谱与分析波段}{0.78}{fig:raw-sic}
\figref{fig:raw-sic}中的大幅反射结构与高波数弱条纹属于不同尺度。若在全谱上直接进行常折射率正弦拟合，强反射带和背景会支配目标函数，所得频率不一定对应透明区中的外延干涉。碳化硅的主分析选择$2000$--$3300\cmiv$，并排除$2300$--$2400\cmiv$的局部异常段；该屏蔽仅作为数据驱动的稳健处理，缺少空白参考谱时不把异常来源确定为某一种气体。改变窗口和屏蔽范围后的结果将在检验部分报告。

\paperfigure{figures/raw_silicon.png}{硅双角度原始反射光谱}{0.78}{fig:raw-si}
硅光谱在低波数处出现较尖锐的谷形，同时包络随波数显著变化，如\figref{fig:raw-si}所示。尖锐条纹是检验高阶反射的重要线索，但材料色散、吸收和背景也会造成非正弦形状。因此，后文将必要条件、曲线形状和模型比较分开讨论，不仅根据图形外观作结论。

\paperfigure{figures/sic_window.png}{碳化硅主分析窗口内的原始弱条纹与局部屏蔽段}{0.78}{fig:sic-window}
图\ref{fig:sic-window}显示，主窗口内的条纹相对于全谱反射平台幅度较小，局部窄带异常却足以影响峰位和包络。最终拟合因此保留缓变背景，并对两角采用同一屏蔽规则。
\subsection{保持测量与加工结果可区分}
原始工作簿作为只读数据保存，并记录SHA-256校验值；由它们生成的CSV采用相同的数值，不预先去趋势或平滑。傅里叶初值与峰位检查使用的辅助平滑结果只用于定位，最终拟合仍面向相应窗口内的实测反射率。每个窗口的选择、有效点数及模型参数均由计算脚本输出。

对本题不执行缺失值插补、聚类剔除或相关性筛选。这里的两个字段分别是自变量和响应量，既不存在待预测样本的特征筛选任务，也没有已确认的缺失观测。这样处理能够保留条纹位置与局部形态，减少加工过程对厚度反演的影响。
